\section{An elliptic family of embedded primary boundaries}
\label{sec:elliptic-boundary}
The universal theorem has a global consequence in a non-isotrivial
family of algebras of Hilbert function $(1,3,3)$. In this section the
parameter tensor is specialized, so its primary decomposition is
proved directly; no compatibility of primary components with arbitrary
base change is presumed.

Let
\begin{equation}\label{eq:elliptic-algebra}
 B_\lambda=\C[x,y,z]/\bigl((x,y,z)^3,
        xy-\lambda z^2,yz-\lambda x^2,zx-\lambda y^2\bigr).
\end{equation}
The elements $1,x,y,z,x^2,y^2,z^2$ are a basis, also over the
polynomial parameter ring. Thus this is a finite free rank-seven
family. Write $V=\langle x,y,z\rangle$ and
$S=\langle x^2,y^2,z^2\rangle$. Work over the nonempty open
\begin{equation}\label{eq:elliptic-parameter-open}
 \Omega=\{\lambda:\lambda\ne0,\quad\lambda^3\ne1,\quad
                  (4\lambda^3-1)^3\ne27\lambda^3\}.
\end{equation}
For example $\lambda=1/2$ belongs to $\Omega$.
Let $X=\Gr(3,V\oplus S)$ and
$U=\{K:\rank(K\to V)\ge2\}\subset X$.
Its determinant divisor $\Delta$ has projection rank exactly two
and is independent of $\lambda$. On $\Pj(V^*)$ put
\begin{equation}\label{eq:elliptic-cubic}
 C_\lambda:\quad F_\lambda(a,b,c)
   =\lambda(a^3+b^3+c^3)+(1-4\lambda^3)abc=0.
\end{equation}
Denote by $E_\lambda\subset\Delta$ its inverse image under
$K\mapsto H=\operatorname{im}(K\to V)$.

\begin{theorem}[Elliptic embedded boundary]
\label{thm:elliptic-boundary}
For every $\lambda\in\Omega$ the curve $C_\lambda$ is smooth and
has genus one. For all $m\ge2$, the full failure scheme $D_\lambda$
of $\Sym^m(\C\oplus K)\to B_\lambda$ on all of $U$ has the
irredundant primary decomposition
\begin{equation}\label{eq:elliptic-primary}
 \mathcal I_{D_\lambda}
     =\mathcal I_\Delta\mathcal I_{E_\lambda}
     =\mathcal I_\Delta\cap\mathcal I_{E_\lambda}^{2}.
\end{equation}
Its only associated points are the generic points of $\Delta$
and $E_\lambda$. The reduced divisor is smooth and has generic
multiplicity one. The embedded support is smooth and normal, with
\begin{equation}\label{eq:elliptic-bundle}
 E_\lambda\simeq
 \operatorname{Tot}_{C_\lambda\times\Pj(S)}
                    \operatorname{Hom}(H,S/\Lambda),
\end{equation}
where $H$ is the tautological hyperplane bundle and $\Lambda$
is the tautological line of $\Pj(S)$. This is a rank-four bundle,
so $\dim E_\lambda=7$ and $\codim_U E_\lambda=2$.

The nilradical has square zero and is the line bundle
$\mathcal O_{E_\lambda}(-\Delta)$, viewed as an
$\mathcal O_{D_\lambda}$-module supported on $E_\lambda$.
More precisely there is an exact sequence
\begin{equation}\label{eq:elliptic-nilradical}
 0\longrightarrow\mathcal O_{E_\lambda}(-\Delta)
 \longrightarrow\mathcal O_{D_\lambda}
 \longrightarrow\mathcal O_\Delta\longrightarrow0.
\end{equation}
For every $K\in\Delta$, the conductor of $\C\oplus K$ in
$B_\lambda$ is exactly $\Lambda=K\cap S$.

The family has genuine moduli: the elliptic curves have nonconstant
$j$-invariant, and distinct $j$-values give nonisomorphic algebras
$B_\lambda$. A smooth projective compactification of $E_\lambda$
is a projective bundle over $C_\lambda\times\Pj^2$ and has
Albanese variety $C_\lambda$.
\end{theorem}
\begin{proof}
A hyperplane with equation $ax+by+cz=0$, on the chart $c\ne0$,
has basis $u=(c,0,-a)$ and $v=(0,c,-b)$. The restriction of the
quadratic multiplication, in the product basis $u^2,uv,v^2$, is
\begin{equation}\label{eq:elliptic-restriction-matrix}
 A_\lambda=
 \begin{pmatrix}
 c^2&-\lambda ac&-2\lambda bc\\
 -2\lambda ac&-\lambda bc&c^2\\
 a^2&ab+\lambda c^2&b^2
 \end{pmatrix},
 \qquad \det A_\lambda=-c^3F_\lambda(a,b,c).
\end{equation}
Expansion along the first row proves the determinant identity.
Cyclic charts give the same equation up to the transition unit of
$\det(\Sym^2H)^*$. Since $\det(\Sym^2H)=(\det H)^3$,
this determinant is a section of $\mathcal O_{\Pj(V^*)}(3)$.
Thus \eqref{eq:elliptic-cubic} defines its global degeneracy curve,
not a component seen only on one chart.

Divide $F_\lambda$ by $\lambda$ and set
\[
 k=\frac{4\lambda^3-1}{3\lambda}.
\]
This gives the Hesse cubic $a^3+b^3+c^3-3kabc=0$.
Its three partial derivatives vanish simultaneously at a projective
point exactly when $k^3=1$: no coordinate of such a point is zero,
and multiplication of $a^2=kbc$, $b^2=kac$, $c^2=kab$ forces
$k^3=1$; the converse follows by solving these equations.
Condition \eqref{eq:elliptic-parameter-open} therefore makes
$C_\lambda$ smooth. It is integral, since distinct components of
a reducible plane cubic intersect and give singularities. Adjunction
gives genus one, and $[1:-1:0]$ is a point on every fibre.
Moreover $A_\lambda$ has rank exactly two along the curve. If
its rank were at most one, all its cofactors, hence all derivatives
of its determinant in the local hyperplane coordinates, would vanish,
contrary to smoothness.

We next check the first-order-spanning hypothesis at every hyperplane,
not only at a general point of the curve. A nonzero functional of $S$
annihilating $\gamma(HV)$ would give a quadratic form in the dual
net vanishing on $HV$. Such a form must be a square of the equation
of $H$, up to a scalar. The dual net consists of
\[
 \alpha(X^2+2\lambda YZ)+\beta(Y^2+2\lambda ZX)
                      +\chi(Z^2+2\lambda XY).
\]
A square $(aX+bY+cZ)^2$ belongs to it only if
\[
 ab=\lambda c^2,\qquad bc=\lambda a^2,\qquad ca=\lambda b^2.
\]
When $\lambda\ne0$, a nonzero solution has all coordinates nonzero.
Their product gives $\lambda^3=1$, excluded on $\Omega$.
It follows that $\gamma(HV)=S$ for every $H$.

Now fix $K\in\Delta$. The block elimination of
\eqref{eq:quadratic-weighted-normal} is valid for this particular
$\gamma$. If $H\notin C_\lambda$, the restriction has rank three,
and the local ideal is $(t)$. If $H\in C_\lambda$, invert a
nonzero two-minor of $A_\lambda$. Eliminating it leaves a one-row
presentation, so the local ideal is
\begin{equation}\label{eq:elliptic-local-ring}
 (t^2,tf),
\end{equation}
where $t$ is a local equation for $\Delta$ and $f$ is the remaining
Schur entry, a unit times $F_\lambda$ in local hyperplane
coordinates. Since $C_\lambda$ is smooth and the map
$\Delta\to\Pj(V^*)$ is smooth, $t,f$ are part of a regular
system of parameters. Off $\Delta$, $M$ is invertible and the
failure scheme is empty. This proves the ideal calculation on all
of $U$ without specializing a universal primary decomposition.

The intrinsic graph description of corank-one planes gives
\eqref{eq:elliptic-bundle}, so $E_\lambda$ is smooth and integral.
Let $P=(t,f)$ locally. Its associated graded ring is a polynomial
ring in two variables over the integral ring of $E_\lambda$.
Consequently $P^2$ is $P$-primary and $(P^2:t)=P$.
It follows that $(t)\cap P^2=tP$, proving
\eqref{eq:elliptic-primary} and its irredundancy. The two intrinsic
ideals glue, and Lemma~\ref{lem:frame-primary-descent} supplies
an alternative frame-to-scheme verification. The ideal
$\mathcal I_\Delta$ is invertible, and therefore
\[
 \mathcal I_\Delta/
          (\mathcal I_\Delta\mathcal I_{E_\lambda})
       \simeq\mathcal I_\Delta\otimes\mathcal O_{E_\lambda}.
\]
This is the kernel in \eqref{eq:elliptic-nilradical}; its square
vanishes because $\mathcal I_\Delta\subset\mathcal I_{E_\lambda}$.
It is nonzero along $E_\lambda$. This proves all nilradical and
multiplicity assertions.

For the conductor, multiplication by a vector $(x,y,z)\in V$ has
matrix
\[
 T(x,y,z)=
 \begin{pmatrix}
 x&\lambda z&\lambda y\\
 \lambda z&y&\lambda x\\
 \lambda y&\lambda x&z
 \end{pmatrix}.
\]
For nonzero $(x,y,z)$ this matrix has rank at least two on $\Omega$.
Indeed its principal two-minors imply, if its rank is at most one,
$xy=\lambda^2z^2$, $yz=\lambda^2x^2$, and $zx=\lambda^2y^2$.
No coordinate can then vanish. The minor in rows $1,2$ and columns
$1,3$ also gives $x^2=\lambda yz$. Combining the latter equality
with $yz=\lambda^2x^2$ gives $\lambda^3=1$, a contradiction.
Thus no nonzero element of $H$ multiplies all of $V$ into a line
$\Lambda$. Proposition~\ref{prop:universal-conductor-flag} gives
$\operatorname{cond}(W)=\Lambda$. Notice also that $W$ is not a
subalgebra on this locus, since $\rank\gamma(\Sym^2H)\ge2$.
Here the conductor means the largest ideal contained in $W$;
it does not presuppose that $W$ is a subalgebra.

Finally the Hesse Weierstrass calculation gives
\begin{equation}\label{eq:elliptic-j}
 j(C_\lambda)
      =\frac{27k^3(k^3+8)^3}{(k^3-1)^3},
 \qquad k=\frac{4\lambda^3-1}{3\lambda}.
\end{equation}
For normalization, substitute $(u_0,u_1)=(1,-k/2)$ in
\cite[Section 2, equations (4)--(5)]{AD09} and use
$j=1728\,4A^3/(4A^3+27B^2)$ for its Weierstrass coefficients.
This yields \eqref{eq:elliptic-j}. It is a nonconstant rational
function of $\lambda$. An isomorphism of local algebras induces
isomorphisms of $\mathfrak m/\mathfrak m^2$ and $\mathfrak m^2$
intertwining $\gamma$. It therefore carries the hyperplane
restriction degeneracy curve projectively to the corresponding
curve. Distinct $j$-values give nonisomorphic algebras.

The compactification
\[
 \overline E_\lambda=
 \Pj_{C_\lambda\times\Pj(S)}
       \bigl(\mathcal O\oplus\operatorname{Hom}(H,S/\Lambda)\bigr)
\]
contains $E_\lambda$ as the affine-bundle chart. It is smooth
and projective. A morphism from a projective space to an abelian
variety is constant; applying this first on its projective-bundle
fibres and then on the $\Pj(S)$ fibres shows that its Albanese
is the Jacobian of $C_\lambda$, identified with $C_\lambda$ by
the specified point. Equivalently, projective-bundle cohomology
gives $h^1(\mathcal O_{\overline E_\lambda})=1$.
This gives a global geometric invariant carried by the embedded
support although the reduced determinant divisor stays fixed.
\end{proof}

\begin{corollary}[A non-isotrivial family of full failure schemes]
\label{cor:elliptic-family}
Over $\Omega$, the construction above gives relative schemes
$\mathcal E\subset\Omega\times\Delta$ and
$\mathcal D\subset\Omega\times U$ with
\[
 \mathcal I_{\mathcal D}
   =\mathcal I_{\Omega\times\Delta}\mathcal I_{\mathcal E},
 \qquad
 0\to\mathcal O_{\mathcal E}(-\Delta)
   \to\mathcal O_{\mathcal D}
   \to\mathcal O_{\Omega\times\Delta}\to0.
\]
The morphisms $\mathcal E\to\Omega$ and $\mathcal D\to\Omega$
are flat; the former is smooth. Each geometric fibre has exactly
the two associated supports in Theorem~\ref{thm:elliptic-boundary}.
\end{corollary}
\begin{proof}
The cubic equation defines a smooth proper family of plane curves
on $\Omega$, since every fibre has nonvanishing relative gradient.
The graph construction is a vector bundle over its product with
$\Pj(S)$, hence $\mathcal E\to\Omega$ is smooth. The local proof
of \eqref{eq:elliptic-local-ring} works over the parameter ring:
its two-minor is invertible and the residual equation $f$ cuts out
that smooth family. It gives the displayed ideal and exact sequence
before specialization. Both outer terms of the sequence are flat
over $\Omega$, so the middle term is flat. Tensoring the sequence
with any residue field is exact, and the fibre proof above gives
the stated primary decomposition. Flatness here is a conclusion of
this specific family, not a general assertion about Fitting schemes.
\end{proof}
